\section{高斯情形下负 KL 散度的解析解}
\label{ap:kl_solution}
变分下界，即需要最大化的目标，包含一个通常可以解析积分的 KL 项。这里给出先验 $\pT(\bz)=\mathcal{N}(0,\bI)$ 和近似后验 $\qPhi(\bz|\bxi)$ 均为高斯分布时的解。令 $J$ 为 $\bz$ 的维数，$\bmu$ 和 $\bsigma$ 分别表示在数据点 $i$ 处计算的变分均值与标准差，$\mu_j$ 和 $\sigma_j$ 表示相应向量的第 $j$ 个分量。以下为简便省略条件 $\bxi$，则有：
\begin{align*}
\int \qPhi(\bz) \log p(\bz) \,\mathrm{d}\bz
&= \int \mathcal{N}(\bz;\bmu,\operatorname{diag}(\bsigma^2)) \log \mathcal{N}(\bz;\bzero,\bI) \,\mathrm{d}\bz \\
&=  - \frac{J}{2} \log (2 \pi) - \frac{1}{2} \sum_{j=1}^J (\mu_j^2 + \sigma_j^2)
\end{align*}
以及：
\begin{align*}
\int \qPhi(\bz) \log \qPhi(\bz) \,\mathrm{d}\bz
&= \int \mathcal{N}(\bz;\bmu,\operatorname{diag}(\bsigma^2)) \log \mathcal{N}(\bz;\bmu,\operatorname{diag}(\bsigma^2)) \,\mathrm{d}\bz \\
&= - \frac{J}{2} \log (2 \pi) - \frac{1}{2} \sum_{j=1}^J ( 1 + \log \sigma^2_j )
\end{align*}
因此：
\begin{align*}
- D_{KL}(\qPhi(\bz) || \pT(\bz)) &= \int \qPhi(\bz) \left(\log \pT(\bz) - \log \qPhi(\bz)\right) \,\mathrm{d}\bz \\
&= \frac{1}{2} \sum_{j=1}^J \left(1 + \log ((\sigma_j)^2) - (\mu_j)^2 - (\sigma_j)^2 \right)
\end{align*}
使用识别模型 $\qPhi(\bz|\bx)$ 时，均值 $\bmu$ 与标准差 $\bsigma$ 就是 $\bx$ 和变分参数 $\bphi$ 的函数，正文给出了相应例子。
